%!TEX program = lualatex
\documentclass[ngerman,headheight=70pt]{scrartcl}
\usepackage[ngerman]{babel}
\usepackage{blindtext}
\usepackage[utf8]{luainputenc}
\usepackage{scrlayer-scrpage}
\usepackage[T1]{fontenc}
\usepackage{fontspec}
\usepackage{multicol}
\usepackage{enumitem}
\usepackage[german=quotes]{csquotes}
\usepackage[bottom=2cm,footskip=8mm]{geometry}
\setmainfont{Source Sans Pro}
\MakeOuterQuote{"}

\setlength{\parindent}{0pt}
\setlength{\parskip}{6pt}

\pagestyle{plain.scrheadings}
\thispagestyle{scrheadings}
\lohead{\vspace*{-1.5cm}\includegraphics[scale=0.4]{../images/stupa_logo_rot.pdf}}
\rohead{\\Präsidium des\\Studierendenparlaments}

\lofoot[Universität Hamburg · Präsidium des Studierendenparlaments\\\hrule Von-Melle-Park 5 · D-20146 Hamburg]{Universität Hamburg · Präsidium des Studierendenparlaments\\\hrule Von-Melle-Park 5 · D-20146 Hamburg}
\cofoot[]{}
\cohead[\thepage]{}

\renewcommand{\thesubsection}{TOP \arabic{subsection}}
\renewcommand{\thesubsubsection}{\arabic{subsubsection}.}
\setcounter{subsection}{-1}

\begin{document}
    UHH · StuPa-Präsidium · Von-Melle-Park 5 · D-20146 Hamburg

    %%% PAGE ONE (EDIT LATER) %%%%
    % Hier den Titel des Protokolls eintragen
    \section*{Protokoll der 2. Sitzung des Studierendenparlaments vom 12. Mai 2016}

    % Protokollanten eintragen
    \textbf{Protokoll: Gunhild Berdal}\\
    \textbf{Ort: S08 in der ehemaligen HWP}\\
    \textbf{Beginn: 18.40 Uhr}\\
    \textbf{Ende: 0.02 Uhr}

    % Wohl die komplizierteste Stelle im Protokoll
    % Als Faustregel sollten nicht mehr als 3 Namen auf einer Zeile stehen
    % Einfach an dieser Vorlage orientieren
    % Im Idealfall ist hier nicht viel zu ändern, da kaum jemand fehlt.
    % Rücktritte müssen natürlich pro Protokoll angepasst werden
    \vspace{0.5cm}
    \begin{tabular}{ll}
        Anwesend: & \\
            RCDS (5 Sitze): & Ramon Weilinger, Antonia Niecke, Ramin Shakiba, \\
                            & Jennifer Maack \\
             CampusGrün (14 Sitze): & Geoffrey Youett, Philipp Droll, Melf Johannsen, \\
                                   & Tahnee Herzig, Mario Moldenhauer, Jim Martens,\\
                                   & Mirzo Ulugbek Khatamov, Armin Günther, Martin Sievert\\
             Bier-Liste (2 Sitze) : & Jakob Pape \\
             WiWi (2 Sitze): & Claas-Friso Hente \\
             Unicorns (5 Sitze): & Johannes Peplow, Annkathrin Löffler, \\
                                 & Andreas Hartkamp, Marielle Hermstrüwer \\
             Liste LINKS (3 Sitze): & Gunhild Berdal, Till Petersen, Sinah Mielich \\
             HWP (2 Sitze): & Samet Gunay, Ajdina Karahasan \\
             MIN (4 Sitze): & \\
             SDS* (3 Sitze): & Mena Winkler, Jacob Petersein, Artur Brückmann \\
             Bart-LISTE (2 Sitze): & Dominic Laumer \\
             LHG (1 Sitz): & Tobias Heisig \\
             harte zeiten (1 Sitz): & Tobias Berking \\
             Jura (1 Sitz): & \\
             AL (2 Sitze): & Karima Schulze, Henri Weber \\
            & \\
        Entschuldigt: & Benjamin Welling (RCDS), Elena Rysikova, Yasemin Günther,\\
                      & Freya Schmitz, Svenja Horn (alle Campus Grün), Karen Martirosian (Bier), \\
                      & Kay Zöllmer (WiWi), Katharina Kucza (Unicorns), Ailina Salten, \\
                      & Lotte Rullkötter, Thea Wahlers (alle MIN), Timo Zeimet (Bart)\\
                      &\\
        Unentschuldigt abwesend: & Laura Franzen (Campus Grün), Jan Detampel (MIN), \\
                                & Johann Baumhoefener (Jura)\\
                                &\\
        Rücktritte: & \\
    \end{tabular}
    \newpage
    %%% Vorgeschlagene Tagesordnung %%%
    \underline{Vorgeschlagene Tagesordnung}
    \begin{enumerate}[label={\textbf{Top \theenumi}},leftmargin=*]
        % folgende Zeile bei Sitzungsfortsetzungen einkommentieren und x durch
        % die Zahl des zuletzt behandelten TOPs ersetzen
        %\setcounter{enumi}{x}
        \item Alle Frauen*-Wahl
            \begin{enumerate}
                \item Bestätigung der Wahlniederschrift (V1617-020)
                \item Bestätigung der Referent*innen (V1617-019)
            \end{enumerate}
        \item Gründung des FSR kath. Theologie (V1617-017) (10 Min.)
        \item Antrag auf VV (V1617-018) (30 Min.)
        \item EM 2016 (V1617-014) (30 Min.)
        \item Verlegung der StuPa-Wahl in das Sommersemester (V1617-015) (30 Min.)
        \item Umgang mit der AfD (V1617-016) (30 Min.)
        \item Verschiedenes
    \end{enumerate}

    \newpage

    % Hier beginnt Teil A. Bei Fortsetzungen kommt hier nur der Teil über
    % die Feststellung der endgültigen Fassung der Tagesordnung.
    % Für sonstige Sitzungen ist dies bereits das Gerüst für die Formalia.
    % Man muss nur noch die protokollierten Inhalte eintragen.

    {\Large \underline{Teil A}}

    \subsection{Formalia}

    \textit{Jim eröffnet die Sitzung um 18.40 Uhr.}

    \subsubsection{Geschäftsbericht Präsidium}

    Das Präsidium hat sich über eine Verbesserung der Öffentlichkeitsarbeit
    verständigt. Insbesondere eine Berichterstattung an FSRe und andere
    Interessierte über Diskussionen und Beschlüsse des StuPas soll versucht
    werden, z.B. über einen Newsletter. Die neue Website wird ständig erweitert.

    Wissenschaftssenatorin Fegebank soll nach dieser Sitzung mit dem Beschluss
    zur VV gegen die Exzellenz-Initiative eingeladen werden.

    Was die Sitzungsplanung betrifft: Am 23. Juni finden gleich drei
    hochschulpolitisch wichtige Termine statt: die Tagung in der
    Erziehungswissenschaft zu Noten,  der APH sowie eine Sitzung des
    Wissenschaftsausschusses. Der Tag ist daher für eine StuPa-Sitzung nicht so
    gut geeignet. Es wird daher überlegt, die Sitzung vom 23. Juni auf den 16.
    Juni vorzuverlegen.

    Daraus würden sich diese Sitzungstermine ab dem 16. Juni ergeben:
    \begin{itemize}
        \item 16. Juni
        \item 30. Juni
        \item 14. Juli (bei Bedarf)
    \end{itemize}

    \subsubsection{Anfragen an das Präsidium}

    Keine

    \subsubsection{Geschäftsbericht AStA}

    \textbf{Philipp (Vorstand)} berichtet, dass der AStA sich konstituiert hat:
    Die Referatszuschnitte sind neu, Organisatorisches hat viel Zeit in Anspruch
    genommen.

    Die Veranstaltung mit Torsten Bultmann (BdWi) war gut: Es wurde insbesondere
    die steigende Abhängigkeit der Hochschulen von Drittmitteln und die
    strukturelle Scheinautonomie der Hochschulen kritisiert.

    In der Frage Akkreditierung von Hochschulen und Studiengängen hat der AStA
    noch keine Linie: Die Vereinheitlichung sei eher positiv zu bewerten, negativ
    hingegen, dass die BA/MA-Studiengänge gefestigt werden.

    \textbf{Franziska (Vorstand)} berichtet von der Demo gegen die Kürzungen in
    Holzwirtschaft, von „Uni für alle statt Exzellenz“. Eine Demo gegen die
    Exzellenzinitiative soll am 14. Juni stattfinden und eine Vollversammlung
    am 2. Juni.

    Sie berichtet weiter kritisch von der Einrichtung von Gebührensatzungen für
    die Studiengänge "Master of International Taxation" (oder besser Master of
    Panama Papers), "Master of legal Studies" (oder auch Troika-Master) und
    "Master der kriminologischen Studis".

    Zur Rahmenprüfungsordnung soll es einen Diskussionstag in der vorlesungsfreien
    Zeit geben.

    Die Studierenden im AS haben eine Initiative für einen Ausschuss für
    internationale Angelegenheiten gestartet.

    \textbf{Timo (Finanzen)} berichtet, dass es erste Überlegungen für den
    Haushalt der VS gibt.

    \textbf{Artur (Soziales, Ökologie, Soziale Bewegung)} zitiert
    Bundespräsident a.D., Horst Köhler, aus dem Jahr 2010, der zurücktreten
    musste, weil er ausgeplaudert hat, dass die Bundeswehr in aller Welt zur
    Sicherung von Rohstoffen und Transportwegen Krieg macht. Dem AStA sei es
    gelungen, die Bundeswehr von der Jobmesse des Stellenwerks auszuladen.

    Der AStA macht eine Veranstaltung mit Mitarbeiter*innen der Charité in
    Berlin, die mehr Personal für die Pflege durchgesetzt haben.

    \textbf{Armin (Soziales, Ökologie, Soziale Bewegung)} kündigt an, dass der
    AStA eine kritische Veranstaltung zu Tierversuchen vorhat und zu Recycling
    in der Stadt arbeitet.

    \textbf{Geoffrey (HoPo, Öffentlichkeitsarbeit, Kultur)} berichtet von viel
    organisatorischem Kram und von der Absicht, besser mit den Kommiliton*innen
    in Kontakt zu kommen.

    \textbf{Kim (HoPo, Öffentlichkeitsarbeit, Kultur)} berichtet, dass der AStA
    die EM problematisieren und kritisieren will.

    \textbf{Janis (HoPo, Öffentlichkeitsarbeit, Kultur)} erzählt, dass die
    Kulturkurse funktionieren, was nicht so gut läuft, wird geprüft.

    \textbf{Mario (Antifaschismus, Antidiskriminierung, Kritische Wissenschaft)}
    berichtet, dass das Referat viele Projekte hat, aber noch keine richtigen
    Konzepte.

    \textbf{Oliver} berichtet aus dem \textbf{Ausschuss gegen Rechts}, dass die
    Lesung "Bücherverbrennung nie wieder" sehr gut war: Gutes Wetter, viel Besuch,
    gute Texte, es wurde viel gelacht. Nicht nur Pazifismus und Antimilitarismus
    und Sozialismus und so, sondern auch Humor war den Nazis ein Dorn im Auge.
    Der Ausschuss plant den "Tag gegen Rechts".

    \subsubsection{Anfragen an den AStA}

    Es gibt Fragen insbesondere zu den Themenbereichen: Ausrichtung der Aktivitäten
    des AStA gegen die EM und gemeinsame Vorhaben des neuen AStA.

    Antworten: Der Kulturbereich eigne sich gut zur Mobilisierung gegen Rechts.
    Die Aktivitäten gegen die Exzellenzinitiative sei das Wichtigste in der
    kommenden Zeit.

    \textit{Zeitansage: 19.30 Uhr}

    \subsubsection{Dringlichkeitsanträge des AStA}

    Keine

    \subsubsection{Aktuelle Stunde (falls entsprechender Antrag vorliegt)}

    Es lag kein Antrag vor.

    \subsubsection{Feststellung der endgültigen Fassung des Teils B der Tagesordnung}

    Antrag Till: TOP 5 "Verlegung der StuPa-Wahl in das Sommersemester" und
    TOP 6 "Umgang mit der AfD" tauschen, damit die Gegen-Rechts-Themen im
    Zusammenhang behandelt werden.

    Mit m:9:0 angenommen.

    Antrag Claas-Friso: TOP 5 zu TOP 4 machen. Mit 2:m:3 abgelehnt.

    Antrag Claas-Friso: TOP 5 zu TOP 1 machen. Mit 8:m:1 abgelehnt.

    Die so geänderte Tagesordnung wird mehrheitlich angenommen.

    %%% BESCHLOSSENE TAGESORDNUNG %%%
    \underline{Endgültige Fassung der Tagesordnung}
    \begin{enumerate}[label={\textbf{Top \theenumi}},leftmargin=*]
        % folgende Zeile bei Sitzungsfortsetzungen einkommentieren und x durch
        % die Zahl des zuletzt behandelten TOPs ersetzen
        %\setcounter{enumi}{x}
        \item Alle Frauen*-Wahl
            \begin{enumerate}
                \item Bestätigung der Wahlniederschrift (V1617-020)
                \item Bestätigung der Referent*innen (V1617-019)
            \end{enumerate}
        \item Gründung des FSR kath. Theologie (V1617-017) (10 Min.)
        \item Antrag auf VV (V1617-018) (30 Min.)
        \item EM 2016 (V1617-014) (30 Min.)
        \item Umgang mit der AfD (V1617-016) (30 Min.)
        \item Verlegung der StuPa-Wahl in das Sommersemester (V1617-015) (30 Min.)
        \item Verschiedenes
    \end{enumerate}

    \subsubsection{Feststellung der Beschlussfähigkeit}

    \textit{Beschlussfähigkeit ist mit 31 von 47 Parlamentarier*Innen gegeben.}

    \subsubsection{Genehmigung der Protokolle der vorangegangenen Sitzungen}

    Till problematisiert die verkürzte Zusammenfassung einzelner Redebeiträge im
    Protokoll des 1. Teils der ersten StuPa-Sitzung, z.B. auf S. 10 "Jakob
    Petersein erwidert...". Das Präsidium ändert die Stelle und das Protokoll
    wird bestätigt.

    Till beantragt die Vertagung des Protokolls vom 2. Teil der ersten
    StuPa-Sitzung mit der Bitte, ähnliche Stellen entsprechend zu ändern -- kein
    Widerspruch.

    \textit{Präsidiumspause bis 20.09 Uhr}

    \vspace{0.5cm}
    {\Large \underline{Teil B}}

    % Teil B beginnt
    % Dieser wird sich für jede Sitzung unterscheiden.
    % Jeder TOP ist eine subsection.
    % Vorhandene Beispiele anfangs entfernen und durch tatsächliche TOPs ersetzen

    % Folgende Zeile bei Sitzungsfortsetzungen einkommentieren und x mit
    % der Zahl des zuletzt behandelten TOPs ersetzen.
    %\setcounter{subsection}{x}
    \subsection{Alle Frauen*-Wahl}

    a) Bestätigung der Wahlniederschrift

    Franziska, Mena und Paula stellen die Wahlniederschrift vor. Es gibt keine
    Fragen. Sie wird mit m:0:2 angenommen.

    b) Bestätigung der Referent*Innen

    Referentinnenvorschlag:
    \begin{itemize}
        \item Melina Toelle (Referentin)
        \item Louisa Escher (Stellvertreterin)
    \end{itemize}

    Melina und Louisa stellen sich vor. Sie legen dar, dass sie vom feministischen
    Standpunkt aus arbeiten: Die gesellschaftliche Ungleichheit werde durch
    Ausbeutung und Machtgefälle im Kapitalismus stabilisiert. Es gelte, patriarchale
    Strukturen aufzubrechen. Konkurrenz verschärfe die Ungleichheit. Aktuelle
    Fragen seien: Wie (sozial) offen ist die Uni? Sexualisierte Gewalt sichtbar
    machen, ein solidarisches Miteinander schaffen. Den "Tag gegen Rechts"
    mitgestalten.

    Es gibt eine engagierte Fragerunde u.a. zu: (Nicht-)Besetzung leerer Stellen
    mit Frauen, zum behaupteten Ausbeutungsverhältnis zwischen Mann und Frau, zu
    positiver Diskriminierung und der Karte für den Campus, auf der Orte
    sexueller Belästigung gekennzeichnet werden, sowie zu Patriotismus, Sexismus
    und (männlicher) Leistungsideologie bei der EM.

    Die Referentinnen betonen, dass sie vor allem die Strukturen der Ausbeutung
    kritisieren, die niemandem gut tun. Wichtig sei die Erkenntnis, dass man nicht
    alleine ist.

    Es wird offen abgestimmt. Die Referentinnen werden mit m:0:9 bestätigt.

    \textit{Zeitansage: 20.50 Uhr}

    \subsection{Gründung des FSR kath. Theologie}

    Geoffrey stellt den Antrag vor.

    Es entsteht eine kritische Diskussion über das (katholische) Verständnis von
    Fachschaftsratsarbeit, Laizismus und Religionsunterricht in den Schulen.
    Einige plädieren dafür, die Gründung des FSRs erstmal zu beschließen und
    dann zu diskutieren. Es wird auch vorgeschlagen, dass jemand vom AStA mit
    den Antragsstellern/Aktiven spricht, wie man FSR-Arbeit machen kann.

    Änderungsantrag: Nicht der "Fachschaftsrat", sondern die "Fachschaft" soll
    eingerichtet werden (in Überschrift und Text). Es gibt keinen Widerspruch
    gegen den Änderungsantrag.

    Der eigentliche Antrag wird mit m:0:3 angenommen.

    \textit{Zeitansage: 21.04 Uhr}

    \subsection{Antrag auf VV}

    Franziska stellt den Antrag vor: Die Kritik an der Exzellenzinitiative nehme
    bundesweit zu, überall sei sie Thema, z.B. bei der Diskussion über den
    Studiengang "Neuropsychologie". Es gehe darum, mit so vielen wie möglich
    darüber zu sprechen, daher die VV.

    Es gibt Fragen vom RCDS: Soll mit der Kampagne nur verhindert werden, dass
    die Exini weitergeführt wird oder soll es um mehr Geld für die Uni gehen?
    Es sei richtig gegen Kürzungen zu sein, aber warum solle man gegen die
    Unterstützung wichtiger Forschung wie z.B. Klimaforschung sein? Die
    Naturwissenschaften bräuchten eben mehr Geld, DESY sei z.B sehr erfolgreich.

    Es wird entgegnet, das Konzept der Exini sei absurd, ihre Fortsetzung sei
    die Verschärfung der Konkurrenz. Die Exzellenz sei selbstreferenziell, ihre
    Kriterien seien Quantität und Verwertungstauglichkeit. Es gehe bei der
    Forderung nach ausreichender Grundfinanzierung der Hochschulen nicht um
    Gleichmacherei, sondern um Bedarf.

    Der Antrag wird mit m:8:0 angenommen.

    \textit{Präsidiumspause bis 21.50 Uhr}

    \subsection{EM 2016}

    Jakob Pape begründet den Antrag der Oppositionslisten "EM 2016 – Public
    Viewing auf dem Campus ermöglichen" und nimmt dabei auch Bezug auf den Antrag
    von Liste LINKS und anderen: Das Verbindende bei Fußball sei wichtig. In
    Israel hätten palästinensische und israelische Mädchen zusammen gespielt.
    Auch hierzulande könnten Flüchtlinge und andere durch Fußball dieselbe
    Sprache sprechen. Die Medien seien auch mal kritisch, z.B. in Brasilien. Man
    könnte Veranstaltungen machen, wo zuerst die Spiele angeguckt werden und dann
    kritisch diskutiert wird.

    Dagegen wird vielfältig Kritik an der EM unter den Gesichtspunkten
    Nationalismus, Rassismus, Antisemitismus, Verdrängung der sozialen Frage und
    Großevent auf dem Campus geäußert.

    Aus der Opposition wird behauptet, einige wollten die EM verbieten, was nicht
    gegen Antisemitismus helfe und nicht zum Frieden führe.

    Es wird erwidert, es gehe nicht ums Verbieten, sondern darum, als VS etwas
    Alternatives zur Aufklärung zu machen.

    \textit{GO-Antrag von Melf: Verlängerung der Sitzung bis nach TOP 7.
    Kein Widerspruch}

    Die Frage kommt auf, wie mit dem Ersetzungsantrag von Liste LINKS
    u.a. umzugehen ist: Ist es ein Änderungsantrag oder werden die beiden Anträge
    alternativ abgestimmt? Es wird sich darauf geeinigt, so wie mit dem Antrag
    des RCDS zur AfD auf der letzten Sitzung zu verfahren -- den Ersetzungsantrag
    zuerst abstimmen und dann den Hauptantrag -- das Präsidium soll die Frage bis
    zum nächsten Mal dann grundsätzlich klären.

    Der Ersetzungsantrag (siehe Anhang) wird mit m:6:3 angenommen.

    Der Gesamtantrag wird mit m:7:0 angenommen.

    \textit{Präsidiumspause bis 23 Uhr}

    \subsection{Umgang mit der AfD}

    Ramin stellt den Antrag "Auseinandersetzung mit der AfD" vor.

    Till spricht gegen den Antrag und jeden einzelnen Absatz und stellt einen
    Änderungsantrag vor.

    Es entsteht eine Debatte über Antifaschismus, Satire und Humor (mit positivem
    Bezug zur Lesung "Bücherverbrennung nie wieder") auf der einen und die "Sorge
    um innere Sicherheit" (Antrag) auf der anderen Seite.

    \textit{Ein GO-Antrag von Melf auf Schließung der Redeliste wird mit 11:9:4
    angenommen.}

    \newpage
    Änderungsantrag LL/harte Zeiten/SDS*:
    \blockquote{Das StuPa bekräftigt, dass der
    Ausschuss gegen Rechts in Zusammenarbeit mit dem AStA einen "Tag gegen Rechts"
    entwickeln, planen und durchführen möge und wird dies nach Kräften unterstützen.
    Insbesondere möge vorbereitet werden, dass sich an diesem Tag mit der AfD
    kritisch-analytisch auseinandergesetzt wird. Auf der Grundlage der Erkenntnisse
    des "Tages gegen Rechts" soll der Ausschuss gegen Rechts eine politische
    Stellungnahme gegen die AfD ausarbeiten, die dem StuPa zur Beschlussfassung
    vorgelegt wird.}

    Änderungsantrag mit m:7:0 angenommen.

    Gesamtabstimmung: Der so geänderte Antrag wird mit m:0:6 angenommen.

    \subsection{Verlegung der StuPa-Wahl in das Sommersemester}

    \textit{Melf beantragt, den TOP auf die nächste Sitzung zu vertagen. Das
    wird mit m:5:3 angenommen.}

    \subsection{Verschiedenes}

    Franziska kündigt die neue Veranstaltung zur Zivilklausel am 6. Juni 2016 an.

    \newpage

    \section*{Ersetzungsantrag von LiLi und anderen zu EM 2016}

    \subsection*{Herren-Fußball-Europameisterschaft 2016 – und wieder
    Großsportveranstaltung}

    Spätestens seit der Auseinandersetzung um Olympische Spiele liegt für die
    meisten Hamburger offen, dass derartige kommerzielle Großevents vor allem
    ein großes Geschäft für Veranstalter, Sponsoren, Ausstatter und Wettbüros
    sind, auf Kosten der großen Mehrheit der Bevölkerung – kulturell begleitet
    von Leistungsideologie, Doping und nationalen Gefühlen.

    Und Medienhype: Keine Zeitung, die nicht mit Torjubel titelt, keine
    Nachrichten, ohne die neusten Informationen aus den Umkleidekabinen. Die
    gnadenlose Suggestion lautet: hältst Du Fußball (oder Handball, oder Laufen,
    oder Schießen, oder was auch immer) nicht für das wichtigste der Welt, bist
    Du ein durchgeknallter Freak. Ein Brief einer offensichtlich konservativen
    Leserin an die FAZ vom 17. Juni 2006 fasste das so zusammen:

    \blockquote[Renate Jansen, Düsseldorf]{Bitte, reduzieren Sie die Berichterstattung über diese duselige WM, ansonsten
    überlege ich mir ernsthaft, für die Zeit der WM eine F.A.Z.-Pause einzulegen.
    Mich interessieren einzelne oder ganze Körperteile von Spielern nicht,
    ebensowenig die Tipergebnisse von Ministern (haben die sonst nichts zu tun?)
    Bieten Sie bitte seriösen Journalismus, decken Sie die Schwächen der Regierung
    egal welcher Farbe auf und machen Sie vernünftige Vorschläge, wie die Karre
    Deutschland wieder aus dem Dreck kommt.}

    Der AStA der Uni Hamburg wird sich an dieser Gehirnwäsche nicht beteiligen,
    mit der den Menschen vorgeschrieben werden soll, was sie zu denken haben:
    dass ihnen ein Sportereignis das wichtigste im Leben sein möge.

    Um eine kulturelle Alternative zu schaffen, organisiert der AStA kurzfristig
    eine Veranstaltungsreihe. Angestrebt wird, rund um die Spieltage der
    Mannschaft aus Deutschland, jeweils einen Referenten aus dem Land der
    gegnerischen Mannschaft zu gewinnen, der von den dortigen sozialen
    Auseinandersetzungen berichtet.

    Darüber hinaus soll der AStA sich weiterhin bemühen, den Breitensport zu
    fördern und in diesem Zusammenhang die Durchführung eines
    AStA-Breitensport-Fussballturniers prüfen.

\end{document}
